\documentclass[level=1]{kidscode}
\newcommand{\lvA}{\faSeedling\ Just starting}
\newcommand{\lvB}{\faShoePrints\ Getting there}
\newcommand{\lvC}{\faCheckCircle\ Got it}
\newcommand{\lvD}{\faRocket\ Expert}
%% "I can" self-rating table: \begin{icantable} \ican{I can ...} \end{icantable}
\newenvironment{icantable}{\par\small\renewcommand{\arraystretch}{1.45}%
  \tabularx{\linewidth}{@{}>{\RaggedRight}X*{4}{>{\centering\arraybackslash}p{1.9cm}}@{}}
  \toprule
  & \lvA & \lvB & \lvC & \lvD\\ \midrule}
  {\bottomrule\endtabularx\par}
\newcommand{\ican}[1]{#1 & \tickbox & \tickbox & \tickbox & \tickbox\\}
\newcommand{\qnum}[1]{\textbf{\color{kcaccent}#1}\quad}
\newcommand{\pts}[1]{\hfill{\small(#1)}}
\begin{document}
\kctitle{Assessments}{Week 1 survey, week 6 checkpoint and week 12 exit assessment}

\begin{tip}[For students: these are not tests to be scared of]
These papers help your teacher see what you already know and what to teach next.
Nobody fails. If you are not sure, write your best guess, or write \enquote{I am not sure yet}:
that is an honest and useful answer. You can ask for a word to be read aloud at any time.
\end{tip}

\begin{teachernote}[For the teacher]
Print each paper only in the week you use it. Answers, scoring and teacher rubrics are in the
Answer Key (\enquote{Assessment answer keys}). Allow extra time, oral answers or a scribe for any
student who needs it: we are checking coding ideas, not handwriting or spelling.
\end{teachernote}

%% =====================================================================
\kcpart{Week 1}{Interest and Comfort Survey}{About 10 minutes. Not graded. There are no wrong answers.}

Nickname (no real names needed): \rule{6cm}{0.4pt}\hfill Date: \rule{3.5cm}{0.4pt}

\section{1. What do I like?}
Check everything you enjoy.
\begin{multicols}{3}
\begin{checklist}
  \item Drawing and painting
  \item Video or board games
  \item Cartoons and animation
  \item Music and dancing
  \item Stories and comics
  \item Puzzles and riddles
  \item Building (blocks, LEGO)
  \item Animals and nature
  \item Space and science
\end{checklist}
\end{multicols}
Something else I love making or doing: \rule{7cm}{0.4pt}

\section{2. How do I feel about coding right now?}
Circle one face.
\begin{center}
\renewcommand{\arraystretch}{1.6}
\begin{tabular}{cccc}
{\Huge\faFrown[regular]} & {\Huge\faMeh[regular]} & {\Huge\faSmile[regular]} & {\Huge\faGrinStars[regular]}\\
Worried & Not sure & Curious & Super excited\\
\end{tabular}
\end{center}

\section{3. Where am I starting?}
Check one box in each row. Most people check \enquote{Just starting} for many rows in week 1. That is exactly right!
\begin{icantable}
\ican{I can use a mouse or trackpad to drag things on the screen.}
\ican{I can type words on a keyboard.}
\ican{I can follow step-by-step instructions in the right order.}
\ican{I can explain what a computer program is.}
\ican{I can keep trying when something does not work the first time.}
\ican{I can work with a partner and take turns.}
\end{icantable}

\section{4. How do I like to learn?}
Check any that are true for you.
\begin{multicols}{2}
\begin{checklist}
  \item I like to see it done first.
  \item I like to try it myself straight away.
  \item I like to read the instructions.
  \item I like someone to read instructions aloud.
  \item Bigger letters on the screen or paper help me.
  \item A quiet spot helps me think.
  \item Short breaks help me focus.
  \item I like to talk ideas through with a partner.
\end{checklist}
\end{multicols}
In pairs, I would like to start as:\quad \tickbox\ the driver (mouse and keyboard)\quad
\tickbox\ the navigator (reads and helps)\quad \tickbox\ I don't mind

\section{5. My dream project}
By the end of Level 1, I would love to make\dots{} (draw it or write it)
\answerspace{4.2cm}

One question I have about coding:
\answerlines{1}

%% =====================================================================
\kcpart{Week 6}{Mid-Level Checkpoint}{About 15 minutes. Scratch reading, putting steps in order, and bug finding.}
\namedateline

\section{Part A: What will the sprite do?}
Read each script. Do not run it: think it through in your head.

\begin{minipage}[t]{0.42\linewidth}
\begin{scratch}[scale=0.75]
\blockinit{when \greenflag clicked}
\blockpen{pen down}
\blockrepeat{repeat \ovalnum{4}}{
  \blockmove{move \ovalnum{100} steps}
  \blockmove{turn \turnright{} \ovalnum{90} degrees}
}
\end{scratch}
\end{minipage}\hfill
\begin{minipage}[t]{0.55\linewidth}
\qnum{A1} What does the sprite draw? Circle one.\pts{1 point}\par\medskip
a line \qquad a square \qquad a triangle \qquad a circle
\end{minipage}

\bigskip\needspace{9\baselineskip}
\begin{minipage}[t]{0.42\linewidth}
\begin{scratch}[scale=0.75]
\blockinit{when \greenflag clicked}
\blockvariable{set \selectmenu{score} to \ovalnum{0}}
\blockrepeat{repeat \ovalnum{3}}{
  \blockvariable{change \selectmenu{score} by \ovalnum{2}}
}
\blocklook{say \ovalvariable{score}}
\end{scratch}
\end{minipage}\hfill
\begin{minipage}[t]{0.55\linewidth}
\qnum{A2} What number does the sprite say?\pts{1 point}\par
\answerlines{1}
\end{minipage}

\bigskip\needspace{9\baselineskip}
\begin{minipage}[t]{0.42\linewidth}
\begin{scratch}[scale=0.75]
\blockinit{when \selectmenu{space} key pressed}
\blocksensing{ask \ovalnum{Do you like cats?} and wait}
\blockifelse{if \booloperator{\ovalsensing{answer} = \ovalnum{yes}} then}
  {\blocklook{say \ovalnum{Me too!} for \ovalnum{2} seconds}}
  {\blocklook{say \ovalnum{OK, dogs then!} for \ovalnum{2} seconds}}
\end{scratch}
\end{minipage}\hfill
\begin{minipage}[t]{0.55\linewidth}
\qnum{A3} When does this script start?\pts{1 point}\par
\answerlines{1}
\qnum{A4} You type \texttt{no}. What does the sprite say?\pts{1 point}\par
\answerlines{1}
\end{minipage}

\section{Part B: Put the steps in order}
\qnum{B1} These steps for planting a seed are mixed up. Number them 1 to 6 so a robot could follow them.\pts{1 point}
\begin{center}\renewcommand{\arraystretch}{1.5}
\begin{tabular}{|p{1.2cm}|p{8cm}|}
\hline
\textbf{Step} & \textbf{Instruction}\\ \hline
 & Drop the seed into the hole.\\ \hline
 & Fill the pot with soil.\\ \hline
 & Water the pot.\\ \hline
 & Make a small hole in the soil with your finger.\\ \hline
 & Put the pot in a sunny place.\\ \hline
 & Cover the seed with soil.\\ \hline
\end{tabular}
\end{center}

\qnum{B2} This dance is long to write. Rewrite it using a \textbf{repeat} so it is shorter.\pts{1 point}\par
\begin{center}
\fbox{\parbox{0.8\linewidth}{\centering clap, stomp, clap, stomp, clap, stomp, clap, stomp, bow}}
\end{center}
\answerlines{2}

\section{Part C: Find the bug}
\begin{minipage}[t]{0.42\linewidth}
\begin{scratch}[scale=0.75]
\blockinit{when \greenflag clicked}
\blockpen{pen down}
\blockrepeat{repeat \ovalnum{3}}{
  \blockmove{move \ovalnum{100} steps}
  \blockmove{turn \turnright{} \ovalnum{90} degrees}
}
\end{scratch}
\end{minipage}\hfill
\begin{minipage}[t]{0.55\linewidth}
\qnum{C1} Leo wants a \textbf{triangle} but gets a shape that is not closed.
Circle the bug in the script and write the fix.\pts{1 point}\par
\answerlines{2}
\end{minipage}

\bigskip\needspace{9\baselineskip}
\begin{minipage}[t]{0.42\linewidth}
{\small\textbf{Cat sprite:}}\par
\begin{scratch}[scale=0.75]
\blockinit{when \greenflag clicked}
\blockevent{broadcast \selectmenu{start game}}
\end{scratch}
\par\smallskip{\small\textbf{Dog sprite:}}\par
\begin{scratch}[scale=0.75]
\blockinit{when I receive \selectmenu{begin}}
\blocklook{say \ovalnum{Let's play!}}
\end{scratch}
\end{minipage}\hfill
\begin{minipage}[t]{0.55\linewidth}
\qnum{C2} Mei clicks the green flag but the dog never says \enquote{Let's play!}. Why not? How would you fix it?\pts{1 point}\par
\answerlines{3}
\end{minipage}

\section{How am I doing?}
Check one box in each row. Be honest: this helps your teacher help you.
\begin{icantable}
\ican{I can put instructions in the right order. (G1)}
\ican{I can use a repeat loop to draw a pattern. (G2)}
\ican{I can make a sprite react with events and if. (G3)}
\ican{I can use a variable and pick random. (G4)}
\ican{I can use messages and make My Blocks. (G5)}
\ican{I can find a bug and fix it. (G6)}
\end{icantable}

%% =====================================================================
\kcpart{Week 12}{Exit Assessment}{About 20 minutes. Read Python, fix bugs, build something small, explain.}
\namedateline

\section{Part A: Predict (paper only)}
\qnum{A1} Draw what this turtle draws. The turtle starts in the middle, facing right.\pts{1 point}\par
\begin{minipage}[t]{0.5\linewidth}
\begin{pycode}
import turtle
t = turtle.Turtle()
for i in range(3):
    t.forward(100)
    t.left(120)
turtle.done()
\end{pycode}
\end{minipage}\hfill
\begin{minipage}[t]{0.46\linewidth}\vspace{2pt}
\answerspace{3.6cm}
\end{minipage}

\qnum{A2} Write exactly what this program prints.\pts{1 point}\par
\begin{minipage}[t]{0.5\linewidth}
\begin{pycode}
for i in range(3):
    print("Hip")
print("Hooray!")
\end{pycode}
\end{minipage}\hfill
\begin{minipage}[t]{0.46\linewidth}\vspace{2pt}
\answerspace{2.8cm}
\end{minipage}

\needspace{10\baselineskip}\qnum{A3} What does this program print? And what would it print if the first line were \texttt{stars = 5}?\pts{2 points}
\begin{pycode}
stars = 7
if stars > 5:
    print("Wow, a starry night!")
else:
    print("Just a few stars.")
\end{pycode}
With \texttt{stars = 7}: \rule{8cm}{0.4pt}\par\medskip
With \texttt{stars = 5}: \rule{8cm}{0.4pt}

\needspace{15\baselineskip}\qnum{A4} How many squares does this program draw? Which one is bigger?\pts{1 point}
\begin{pycode}
import turtle
t = turtle.Turtle()

def square(size):
    for i in range(4):
        t.forward(size)
        t.left(90)

square(50)
square(100)
turtle.done()
\end{pycode}
\answerlines{2}

\section{Part B: Fix the bugs (paper only)}
\needspace{13\baselineskip}\qnum{B1} Omar's program should ask for a number of sides and draw that shape.
It has \textbf{3 bugs}. Circle each one and write the fixed line next to it.\pts{3 points}
\begin{pycode}
import turtle
t = turtle.Turtle()
sides = input("How many sides? ")
for i in range(sides)
    t.forwad(80)
    t.left(360 / sides)
turtle.done()
\end{pycode}
\answerlines{3}

\needspace{15\baselineskip}\qnum{B2} Zoe runs this program. No error appears, but nothing is drawn. Why? Write the line she needs to add.\pts{1 point}
\begin{pycode}
import turtle
t = turtle.Turtle()

def flower(color):
    t.color(color)
    for i in range(6):
        t.circle(40)
        t.left(60)

turtle.done()
\end{pycode}
\answerlines{2}

\section{Part C: Build it (at the computer, about 7 minutes)}
\qnum{C1} Open Thonny and start a new file. Write a program that:\pts{3 points}
\begin{checklist}
  \item asks the viewer \texttt{How many circles?}
  \item draws that many circles,
  \item puts each circle in a random place.
\end{checklist}
Extra (if you have time): put the circle-drawing in your own function, or give the circles colors.
Save it as \texttt{exit\_circles.py} and show your teacher.

\section{Part D: Explain it}
\needspace{8\baselineskip}\qnum{D1} Look at the \texttt{square(size)} program in A4. In your own words: what does \texttt{size} do,
and why is making a function like \texttt{square} useful? (You can say your answer to your teacher instead of writing it.)\pts{2 points}
\answerlines{4}

\section{How am I doing?}
Check one box in each row. Look back at your week 1 survey: how far have you come?
\begin{icantable}
\ican{I can read Python code and predict what it does. (G9)}
\ican{I can use a \texttt{for} loop to repeat things. (G2)}
\ican{I can use \texttt{input}, variables and \texttt{if}. (G3, G4)}
\ican{I can use \texttt{random} to add a surprise. (G4)}
\ican{I can make and use a function with \texttt{def}. (G5)}
\ican{I can find and fix bugs. (G6)}
\ican{I can explain my code to someone else. (G7)}
\end{icantable}

\end{document}
